\chapter{Chapter 10 Refresher: LLM Systems --- Retrieval, Tool Use, and Evaluation Harnesses --- Retrieval Metrics and Evaluation Arithmetic}

\section*{Setting the Scene}
Retrieval-augmented answering concatenates two stochastic modules: recall relevant evidence (set-valued event), then condition generation on that evidence (Preface Step~2 again). One scalar cannot isolate which module failed.

\paragraph{Step A --- Set-valued relevance.}
\textbf{Problem.} Truth may admit multiple supporting documents.
\textbf{Insight.} Precision/recall generalise Preface Step~1 counting logic to partial overlaps between retrieved and gold sets.
\textbf{Lineage.} Classical IR (Manning et al.); neural hybrids layer embedding geometry from Chapter~4.

\paragraph{Step B --- Conditional generation quality.}
\textbf{Problem.} Good retriever + weak synthesiser still ships nonsense.
\textbf{Insight.} Evaluate \(\mathcal{L}\) or task metrics \emph{given} oracle contexts vs live contexts---difference isolates grounding gaps.

\paragraph{Step C --- Faithfulness as constrained decoding.}
\textbf{Problem.} Fluent generators invent citations.
\textbf{Insight.} Faithfulness metrics approximate hard logical entailment with softer overlap checks---engineering surrogate for impossible full semantics.

Decomposition honors Preface Step~5 linearity instinct: diagnose components before averaging them blindly.

\section*{Notation Introduced in This Chapter}
\begin{itemize}[leftmargin=1.5em]
  \item $k$: retrieval cutoff.
  \item $R_k$: top-$k$ retrieved set.
  \item $G$: gold relevant set.
  \item Precision@$k$, Recall@$k$: standard ranking metrics.
  \item $q$: query embedding/vector.
  \item $d$: document/chunk embedding/vector.
\end{itemize}

\section*{What You Need To Recall Precisely}
\begin{itemize}[leftmargin=1.5em]
  \item \textbf{Ranking metrics.}
  Let \(R_k\) be top-$k$ retrieved IDs and \(G\) gold relevant IDs with $|G|$ possibly larger than $k$.
  \[
    \mathrm{Prec@}k=\frac{|R_k \cap G|}{k},\qquad
    \mathrm{Rec@}k=\frac{|R_k \cap G|}{|G|}.
  \]
  Precision stresses cleanliness of short rankings; recall stresses coverage of all positives---trade-offs arise whenever $|G|$ is large relative to $k$. Mean reciprocal rank (MRR) averages \(1/\text{rank of first relevant}\) across queries---rewarding systems that surface \emph{one} correct doc quickly rather than maximising exhaustive recall.
  \item \textbf{Sparse vs dense retrieval:} lexical overlap and embedding similarity capture different relevance modes.
  \item \textbf{Fusion methods:} rank fusion stabilizes retrieval when signals disagree.
  \item \textbf{Evaluation decomposition:} retrieval-stage errors and generation-stage errors require separate diagnostics.
  \item \textbf{Grounding metrics:} answer support must be checked against retrieved evidence, not linguistic plausibility.
\end{itemize}

\section*{How These Results Build on Each Other}
\begin{enumerate}[leftmargin=1.5em]
  \item Define relevance and evidence assumptions.
  \item Measure retrieval coverage/precision.
  \item Measure generation quality conditioned on retrieved context.
  \item Measure grounding/faithfulness of final claims.
  \item Localize failure stage before proposing fix.
\end{enumerate}

\section*{Reminder Glossary (Term by Term)}
\begin{itemize}[leftmargin=1.5em]
  \item \textbf{Gold relevance set:} trusted set of truly relevant items.
  \item \textbf{Precision@k:} fraction of top-$k$ retrieved items that are relevant.
  \item \textbf{Recall@k:} fraction of relevant items captured in top-$k$.
  \item \textbf{Reranker:} second-stage scoring model for refined ranking.
  \item \textbf{Faithfulness:} proportion of answer claims supported by context.
\end{itemize}

\section*{Worked Mini Example (Precision and Recall@k)}
Suppose the gold relevant set has 4 items. Top-5 retrieval returns 3 relevant items.
\[
\mathrm{Precision@5}=\frac{3}{5}=0.6,\qquad
\mathrm{Recall@5}=\frac{3}{4}=0.75.
\]
Interpretation: ranking quality and evidence coverage are both incomplete, but in different ways.

\section*{Worked Example (Failure Localization)}
Pipeline A improves final answer score from 68 to 72.  
But decomposition reveals:
\begin{itemize}[leftmargin=1.5em]
  \item retrieval recall@5 dropped from 0.81 to 0.70,
  \item generation conditioned on gold context improved substantially.
\end{itemize}
Interpretation: generator improved, retriever regressed.  
Action: fix retrieval stage rather than over-tuning generation prompts.

\section*{Intuition Checkpoints}
\begin{itemize}[leftmargin=1.5em]
  \item Better answer fluency can hide retrieval failures.
  \item High retrieval recall can coexist with poor final answers if synthesis is weak.
  \item Metric definitions require explicit edge-case conventions.
\end{itemize}

\section*{Further Refresh}
\begin{itemize}[leftmargin=1.5em]
  \item Embedding geometry: see Chapter 4 refresher.
  \item Calibration/uncertainty interpretation: see Chapter 1 refresher.
  \item Systems throughput constraints: see Chapter 7 refresher.
  \item \textbf{Manning et al., \emph{IIR}} --- BM25, ranking metrics, IR evaluation.
  \item \textbf{FAISS/HNSW papers} --- ANN retrieval behavior and tradeoffs.
\end{itemize}
